% Options for packages loaded elsewhere
\PassOptionsToPackage{unicode}{hyperref}
\PassOptionsToPackage{hyphens}{url}
\PassOptionsToPackage{dvipsnames,svgnames,x11names}{xcolor}
\documentclass[
  10pt,
]{article}
\usepackage{xcolor}
\usepackage{amsmath,amssymb}
\setcounter{secnumdepth}{-\maxdimen} % remove section numbering
\usepackage{iftex}
\ifPDFTeX
  \usepackage[T1]{fontenc}
  \usepackage[utf8]{inputenc}
  \usepackage{textcomp} % provide euro and other symbols
\else % if luatex or xetex
  \usepackage{unicode-math} % this also loads fontspec
  \defaultfontfeatures{Scale=MatchLowercase}
  \defaultfontfeatures[\rmfamily]{Ligatures=TeX,Scale=1}
\fi
\usepackage{lmodern}
\ifPDFTeX\else
  % xetex/luatex font selection
\fi
% Use upquote if available, for straight quotes in verbatim environments
\IfFileExists{upquote.sty}{\usepackage{upquote}}{}
\IfFileExists{microtype.sty}{% use microtype if available
  \usepackage[]{microtype}
  \UseMicrotypeSet[protrusion]{basicmath} % disable protrusion for tt fonts
}{}
\makeatletter
\@ifundefined{KOMAClassName}{% if non-KOMA class
  \IfFileExists{parskip.sty}{%
    \usepackage{parskip}
  }{% else
    \setlength{\parindent}{0pt}
    \setlength{\parskip}{6pt plus 2pt minus 1pt}}
}{% if KOMA class
  \KOMAoptions{parskip=half}}
\makeatother
\usepackage{color}
\usepackage{fancyvrb}
\newcommand{\VerbBar}{|}
\newcommand{\VERB}{\Verb[commandchars=\\\{\}]}
\DefineVerbatimEnvironment{Highlighting}{Verbatim}{commandchars=\\\{\}}
% Add ',fontsize=\small' for more characters per line
\newenvironment{Shaded}{}{}
\newcommand{\AlertTok}[1]{\textcolor[rgb]{1.00,0.00,0.00}{\textbf{#1}}}
\newcommand{\AnnotationTok}[1]{\textcolor[rgb]{0.38,0.63,0.69}{\textbf{\textit{#1}}}}
\newcommand{\AttributeTok}[1]{\textcolor[rgb]{0.49,0.56,0.16}{#1}}
\newcommand{\BaseNTok}[1]{\textcolor[rgb]{0.25,0.63,0.44}{#1}}
\newcommand{\BuiltInTok}[1]{\textcolor[rgb]{0.00,0.50,0.00}{#1}}
\newcommand{\CharTok}[1]{\textcolor[rgb]{0.25,0.44,0.63}{#1}}
\newcommand{\CommentTok}[1]{\textcolor[rgb]{0.38,0.63,0.69}{\textit{#1}}}
\newcommand{\CommentVarTok}[1]{\textcolor[rgb]{0.38,0.63,0.69}{\textbf{\textit{#1}}}}
\newcommand{\ConstantTok}[1]{\textcolor[rgb]{0.53,0.00,0.00}{#1}}
\newcommand{\ControlFlowTok}[1]{\textcolor[rgb]{0.00,0.44,0.13}{\textbf{#1}}}
\newcommand{\DataTypeTok}[1]{\textcolor[rgb]{0.56,0.13,0.00}{#1}}
\newcommand{\DecValTok}[1]{\textcolor[rgb]{0.25,0.63,0.44}{#1}}
\newcommand{\DocumentationTok}[1]{\textcolor[rgb]{0.73,0.13,0.13}{\textit{#1}}}
\newcommand{\ErrorTok}[1]{\textcolor[rgb]{1.00,0.00,0.00}{\textbf{#1}}}
\newcommand{\ExtensionTok}[1]{#1}
\newcommand{\FloatTok}[1]{\textcolor[rgb]{0.25,0.63,0.44}{#1}}
\newcommand{\FunctionTok}[1]{\textcolor[rgb]{0.02,0.16,0.49}{#1}}
\newcommand{\ImportTok}[1]{\textcolor[rgb]{0.00,0.50,0.00}{\textbf{#1}}}
\newcommand{\InformationTok}[1]{\textcolor[rgb]{0.38,0.63,0.69}{\textbf{\textit{#1}}}}
\newcommand{\KeywordTok}[1]{\textcolor[rgb]{0.00,0.44,0.13}{\textbf{#1}}}
\newcommand{\NormalTok}[1]{#1}
\newcommand{\OperatorTok}[1]{\textcolor[rgb]{0.40,0.40,0.40}{#1}}
\newcommand{\OtherTok}[1]{\textcolor[rgb]{0.00,0.44,0.13}{#1}}
\newcommand{\PreprocessorTok}[1]{\textcolor[rgb]{0.74,0.48,0.00}{#1}}
\newcommand{\RegionMarkerTok}[1]{#1}
\newcommand{\SpecialCharTok}[1]{\textcolor[rgb]{0.25,0.44,0.63}{#1}}
\newcommand{\SpecialStringTok}[1]{\textcolor[rgb]{0.73,0.40,0.53}{#1}}
\newcommand{\StringTok}[1]{\textcolor[rgb]{0.25,0.44,0.63}{#1}}
\newcommand{\VariableTok}[1]{\textcolor[rgb]{0.10,0.09,0.49}{#1}}
\newcommand{\VerbatimStringTok}[1]{\textcolor[rgb]{0.25,0.44,0.63}{#1}}
\newcommand{\WarningTok}[1]{\textcolor[rgb]{0.38,0.63,0.69}{\textbf{\textit{#1}}}}
\usepackage{longtable,booktabs,array}
\usepackage{caption}
\captionsetup[table]{skip=6pt}
\newcounter{none} % for unnumbered tables
\usepackage{calc} % for calculating minipage widths
% Correct order of tables after \paragraph or \subparagraph
\usepackage{etoolbox}
\makeatletter
\patchcmd\longtable{\par}{\if@noskipsec\mbox{}\fi\par}{}{}
\makeatother
% Allow footnotes in longtable head/foot
\IfFileExists{footnotehyper.sty}{\usepackage{footnotehyper}}{\usepackage{footnote}}
\makesavenoteenv{longtable}
\setlength{\emergencystretch}{3em} % prevent overfull lines
\providecommand{\tightlist}{%
  \setlength{\itemsep}{0pt}\setlength{\parskip}{0pt}}
\usepackage[letterpaper,margin=0.75in]{geometry}
\usepackage{booktabs,longtable,array,multirow,colortbl}
\usepackage{fancyhdr}
\usepackage{seqsplit}
\usepackage{etoolbox}
\definecolor{muted}{HTML}{57606A}
% Wide tables: shrink the font rather than run off the page.
\let\oldlongtable\longtable
\def\longtable{\scriptsize\oldlongtable}
% Long monospace names (run slugs, script paths, config keys) may break
% anywhere, in tables and paragraphs alike. Robust, so section titles that
% carry code are written to the .toc/.aux unexpanded instead of crashing.
\AtBeginDocument{\DeclareRobustCommand{\texttt}[1]{{\ttfamily\seqsplit{#1}}}}
\setlength{\tabcolsep}{4pt}
\renewcommand{\arraystretch}{1.12}
\setlength{\parindent}{0pt}
\setlength{\parskip}{5pt}
\setlength{\emergencystretch}{3em}
\pagestyle{fancy}\fancyhf{}
\renewcommand{\headrulewidth}{0pt}
\fancyfoot[L]{\footnotesize\color{muted}CLAUDE.md -- EmbodiedMAE-4M, yongyun branch}
\fancyfoot[R]{\footnotesize\color{muted}\thepage}
\usepackage{bookmark}
\IfFileExists{xurl.sty}{\usepackage{xurl}}{} % add URL line breaks if available
\urlstyle{same}
% fallback for those not using the hyperref driver hyperxmp:
\makeatletter
\@ifundefined{xmpquote}{\newcommand{\xmpquote}[1]{#1}}{}
\makeatother
\hypersetup{
  colorlinks=true,
  linkcolor={RoyalBlue},
  filecolor={Maroon},
  citecolor={Blue},
  urlcolor={RoyalBlue},
  pdfcreator={LaTeX via pandoc}}

\title{CLAUDE.md — EmbodiedMAE-4M repository guide and experiment log}
\author{}
\date{2026-09-28, \texttt{yongyun} branch}

\begin{document}
\maketitle

{
\setcounter{tocdepth}{3}
\tableofcontents
}
\section{CLAUDE.md}\label{claude.md}

This file provides guidance to Claude Code (claude.ai/code) when working
with code in this repository.

\subsection{What this is}\label{what-this-is}

PyTorch implementation of \textbf{EmbodiedMAE} --- a multi-modal Masked
Autoencoder pre-trained on synthetic Sorghum plant data. Two model
variants share the same encoder/decoder skeleton:

\begin{itemize}
\tightlist
\item
  \textbf{3M (\texttt{embodied\_mae.py})} --- RGB + Depth + Point Cloud
\item
  \textbf{4M (\texttt{embodied\_mae\_4m.py})} --- adds a parametric
  ``spline/text'' modality describing the plant's procedural generation
  parameters
\end{itemize}

Both use the same Dirichlet-allocation token masking, transformer
encoder, and per-modality decoder heads. The 4M variant is the active
line of work; 3M is kept as the baseline it was derived from.

\subsection{Repo layout}\label{repo-layout}

Configs and batch scripts were consolidated out of the repo root.
\textbf{There are no \texttt{config.yaml} / \texttt{config\_4m.yaml} at
the root any more} --- everything lives in:

\begin{itemize}
\tightlist
\item
  \texttt{configs/} --- every YAML. \texttt{config.yaml} (3M),
  \texttt{config\_4m.yaml} (4M), \texttt{config\_4m\_pretrain15k*.yaml}
  (pretrain), \texttt{config\_4m\_distill\_*.yaml} (cross-modal
  distillation), \texttt{config\_e2\_*.yaml} (the E2 modality ablation),
  \texttt{config\_e3\_*.yaml} (data scaling) and
  \texttt{config\_e4\_*.yaml} (model scaling).
\item
  \texttt{slurm/} --- every \texttt{.sbatch} / launcher script.
  \texttt{e2\_arm*.sbatch} launch E2 arms;
  \texttt{scale\_arm\_blackwell.sbatch} launches any E3 or E4 arm by
  slug
  (\texttt{sbatch\ -\/-job-name=e3\_1k\ slurm/scale\_arm\_blackwell.sbatch\ e3\_1k}),
  deriving \texttt{configs/config\_\textless{}slug\textgreater{}.yaml}
  and \texttt{outputs/\textless{}slug\textgreater{}/}.
\item
  \texttt{data\_split/} --- the train/val/test split tooling
  (\texttt{make\_split.py}, \texttt{move\_split.py},
  \texttt{reshuffle.py}).
\item
  \texttt{eval/} --- evaluation and analysis
  (\texttt{eval/analyze\_e2.py},
  \texttt{eval/eval\_views\_one\_plant.py},
  \texttt{eval/vis\_pc\_unpredicted.py},
  \texttt{eval/eval\_warmstart.py}, \ldots).
\item
  \texttt{figures/} --- anything that renders a figure, report or deck
  (\texttt{figures/make\_results\_pptx.py}, \texttt{plot\_*.py},
  \texttt{gen\_*.py}, \texttt{regen\_*.py}).
\item
  \texttt{export/} --- dumps and artefact builders
  (\texttt{export/dump\_param\_examples.py}, \texttt{export\_pc\_*.py},
  \texttt{build\_*.py}).
\item
  \texttt{sweeps/} --- masking and visualisation sweeps.
\item
  \texttt{tools/} --- one-off data-prep and diagnostic scripts.
\item
  \textbf{The repo root holds only what other code imports}: the two
  models, the two datasets, the training entry points,
  \texttt{validate*.py} and \texttt{utils.py}. That is the rule --- if a
  file at the root is imported by nothing, it belongs in one of the
  folders above.
\end{itemize}

Scripts in those folders put the repo root on \texttt{sys.path}
themselves (a four-line shim at the top), so
\texttt{python\ eval/analyze\_e2.py} works from the repo root without
\texttt{PYTHONPATH} set. Keep the shim when adding a script there.

If a command or script still refers to a root-level
\texttt{config\_4m.yaml}, it is stale --- the path is
\texttt{configs/config\_4m.yaml}.

\subsection{Experiment programme --- what is done and what is
left}\label{experiment-programme-what-is-done-and-what-is-left}

The programme is the E1--E10 matrix in the CVPR 2027 plan (the Google
Doc is the source of truth for scope; this section is the source of
truth for \emph{state}). Paper deadline \textbf{Nov 13 2026}, internal
results freeze \textbf{Oct 24 2026}.

\textbf{Status as of 2026-09-20.} Run state goes stale fast --- the
table records which experiments have been \emph{built}, which is
durable. For live progress use \texttt{squeue\ -u\ \$USER},
\texttt{outputs/\textless{}run\textgreater{}/training\_history.json},
and \texttt{outputs/\textless{}run\textgreater{}/config.json} (which
records what a run actually used).

{\def\LTcaptype{none} % do not increment counter
\begin{longtable}[]{@{}
  >{\raggedright\arraybackslash}p{(\linewidth - 6\tabcolsep) * \real{0.2500}}
  >{\raggedright\arraybackslash}p{(\linewidth - 6\tabcolsep) * \real{0.2500}}
  >{\raggedright\arraybackslash}p{(\linewidth - 6\tabcolsep) * \real{0.2500}}
  >{\raggedright\arraybackslash}p{(\linewidth - 6\tabcolsep) * \real{0.2500}}@{}}
\toprule\noalign{}
\begin{minipage}[b]{\linewidth}\raggedright
\end{minipage} & \begin{minipage}[b]{\linewidth}\raggedright
experiment
\end{minipage} & \begin{minipage}[b]{\linewidth}\raggedright
state
\end{minipage} & \begin{minipage}[b]{\linewidth}\raggedright
runs
\end{minipage} \\
\midrule\noalign{}
\endhead
\bottomrule\noalign{}
\endlastfoot
E1 & headline pretrain & \textbf{done} &
\texttt{4m\_pretrain\_15k\_v2\_depthfix\_qal}, 1000/1000 ep, global
batch 256 \\
E2 & modality value-add & \textbf{done} & all four arms 600/600; results
in \texttt{reports/RESULTS\_DECK\_2026-09-20.md} \\
E3 & data scaling & \textbf{2 of 3} & \texttt{e3\_3k} yes
\texttt{e3\_10k} yes; \texttt{e3\_1k} resuming from epoch 5654/6169 \\
E4 & model scaling & \textbf{1 of 2} & \texttt{e4\_small} yes;
\texttt{e4\_large} resuming from epoch 250/600 \\
E5 & view regime & \textbf{not built} & --- \\
E6 & masking / noise & \textbf{600-epoch pass done; 1 000-epoch
continuation pending} (\texttt{yongyun} branch) & structured masking
trades clean PC accuracy for robustness: worse clean F1, near-zero
degradation under corruption --- see
\hyperref[e6--e7-on-the-yongyun-branch]{E6 / E7 on the \texttt{yongyun}
branch} \\
E7 & loss study & \textbf{first pass done} (\texttt{yongyun} branch) &
QAL \textgreater{} Chamfer \textgreater{} Sinkhorn-only on F1; Sinkhorn
trades precision for recall; \texttt{model.loss\_name} now also takes
\texttt{sinkhorn} --- see below \\
E8 & baselines & \textbf{not built} & --- \\
E9 & latent analysis & \textbf{not built} & --- \\
E10 & real-data OOD & \textbf{partial} & \texttt{OOD\_EVAL\_rgb2pc.md}
and the \texttt{eval\_rgb2pc\_*.py} scripts \\
\end{longtable}
}

\textbf{On the two resuming arms.} \texttt{e3\_1k} and
\texttt{e4\_large} both died at exactly \texttt{2026-09-19T15:51:27} on
\texttt{nova26-gpu-2} --- same second, same node, no traceback, host RSS
far under request. That is a node-level event on the preemptible
\texttt{scavenger} partition, not a fault in either run, and the
launcher's auto-resume bounds the loss at \texttt{save\_freq} epochs.
\texttt{e3\_1k}'s \texttt{training\_history.json} was left truncated
mid-write; its val series is recoverable from the run log in
\texttt{logs/}.

\textbf{On waiting for GPUs.} When these jobs sit \texttt{PENDING} for
days, check whether it is your request before shrinking it: a 1-GPU /
1-CPU / 1 GB / 10-minute test job placed at the \emph{same} time as the
full 2-GPU / 320 GB request, which means the nodes are held by a
reservation and no amount of trimming helps.

\subsubsection{Two gaps that block the headline
claim}\label{two-gaps-that-block-the-headline-claim}

\textbf{1. There is no downstream linear probe, and the plan says that
is the metric.} Locked decision 6.4 makes the value-add metric a linear
probe on height, leaf angle, leaf count and biomass --- explicitly
\emph{not} reconstruction loss. \texttt{eval/analyze\_e2.py} compares
arms on \texttt{val\_pc\_chamfer}, which is the correct arm-invariant
\emph{monitoring} signal and is not what 6.4 asks for. So E2, E3 and E4
will all finish and produce chamfer curves with no probe number
attached.

Three of the four targets are already in \texttt{features.csv} at the
split root (15 000 rows keyed by \texttt{plant}, which joins to
\texttt{Sorghum\_\textless{}plant\textgreater{}\_\textless{}view\textgreater{}}):
\texttt{stem\_length} \(\to\) height, \texttt{n\_leaves} \(\to\) leaf
count, and leaf angle is either \texttt{roll\_mean} (twist) or
\texttt{branch\_mean} (insertion) --- the spline YAMLs carry
\texttt{roll\_angle} and \texttt{branching\_angle} as distinct fields
and the plan does not say which it means. \textbf{Biomass is in neither
\texttt{features.csv} nor the spline params} and has to be derived
(cheapest proxy: \texttt{n\_leaves\ ×\ leaf\_len\_mean}).

Build and validate this against an existing checkpoint \emph{before} the
arms finish --- if the probe is broken or these targets carry no linear
signal, that is much cheaper to learn now than after the runs are
unrepeatable.

\textbf{2. No baselines exist (E8), and the plan ranks that the \#1
reject risk.} Nothing in the repo addresses it, and every baseline is
itself a training run that has to fit before the Oct 24 freeze, so the
\emph{decision} about what to compare against is more time-critical than
the runs.

\subsubsection{On E6 and E7 being run
elsewhere}\label{on-e6-and-e7-being-run-elsewhere}

Masking/noise and the loss study are a collaborator's, so nothing in
this repo should schedule them. Their results still have to land in the
same table as everything above, which means the comparison only holds if
they match this programme on the three things that silently break it:
the \textbf{same} 70/15/15 split of \texttt{Sorghum\_15K} (seed 42), the
\textbf{same} global batch of 32, and a budget stated in
\textbf{optimizer steps} rather than epochs. Confirm those three before
their numbers are merged, not after ---
\texttt{outputs/\textless{}run\textgreater{}/config.json} records all
three for any run that used this codebase.

That leaves \textbf{E5, E8, E9 and E10 plus the linear probe} as the
work owned here.

\subsubsection{\texorpdfstring{E6 / E7 on the \texttt{yongyun}
branch}{E6 / E7 on the yongyun branch}}\label{e6-e7-on-the-yongyun-branch}

Results from the collaborator side, 2026-09-21. Every run below is
\textbf{4M base, recipe-B occlusion + structured masking, 3 000 plants,
one view per plant per epoch, 30 000 optimizer steps} (94 steps/epoch
\(\to\) 320 epochs; best checkpoints at epochs 300--320, i.e.~still
improving when the cosine schedule ends), evaluated on a fixed 225-view
clean val set and, for the occluded columns, under one reference
corruption shared by all runs (\texttt{occlusion.eval\_occlusion}, 1×
noise).

\subsubsection{The 600-epoch runs (2026-09-26) --- these are the
results}\label{the-600-epoch-runs-2026-09-26-these-are-the-results}

Twelve runs, \texttt{train\_4m\_paper.sbatch} +
\texttt{config\_4m\_paper\_*.yaml}: \textbf{56 400 optimizer steps = 600
epochs} at 94 steps/epoch (3 000 plants, one view per plant per epoch),
one cosine from scratch with 19 warm-up epochs, global batch 32 (8 per
GPU × 4 L40S), fixed 225-view clean val. Arms that carry a claim run at
\textbf{seeds 1 and 2}, so every gap can be read against the run-to-run
spread.

600 epochs is enough here: PC F1@0.03 climbs from 0.60 to 0.73 between
epochs 200 and 250 and is flat over the last four validation points
(e.g.~0.733, 0.733, 0.735, 0.734), so arms are compared after the
transition, not before it. It is, however, \textbf{a short schedule next
to the programme's 197 400 steps} and has to be labelled that way
wherever these rows are reported. An earlier attempt at the full 197
400-step budget was stopped at \textasciitilde580 epochs and is kept in
\texttt{outputs/\_aborted\_2100budget/}; its numbers are not comparable
because its cosine was still at 85 \% of peak LR there.

\paragraph{E6 --- structured masking: worse clean reconstruction, better
robustness}\label{e6-structured-masking-worse-clean-reconstruction-better-robustness}

{\def\LTcaptype{none} % do not increment counter
\begin{longtable}[]{@{}
  >{\raggedright\arraybackslash}p{(\linewidth - 12\tabcolsep) * \real{0.1429}}
  >{\raggedright\arraybackslash}p{(\linewidth - 12\tabcolsep) * \real{0.1429}}
  >{\raggedright\arraybackslash}p{(\linewidth - 12\tabcolsep) * \real{0.1429}}
  >{\raggedright\arraybackslash}p{(\linewidth - 12\tabcolsep) * \real{0.1429}}
  >{\raggedright\arraybackslash}p{(\linewidth - 12\tabcolsep) * \real{0.1429}}
  >{\raggedright\arraybackslash}p{(\linewidth - 12\tabcolsep) * \real{0.1429}}
  >{\raggedright\arraybackslash}p{(\linewidth - 12\tabcolsep) * \real{0.1429}}@{}}
\toprule\noalign{}
\begin{minipage}[b]{\linewidth}\raggedright
arm
\end{minipage} & \begin{minipage}[b]{\linewidth}\raggedright
token masking
\end{minipage} & \begin{minipage}[b]{\linewidth}\raggedright
F1@0.03 s1 / s2
\end{minipage} & \begin{minipage}[b]{\linewidth}\raggedright
recall@0.03 s1 / s2
\end{minipage} & \begin{minipage}[b]{\linewidth}\raggedright
F1@0.01 s1 / s2
\end{minipage} & \begin{minipage}[b]{\linewidth}\raggedright
PC Chamfer s1 / s2
\end{minipage} & \begin{minipage}[b]{\linewidth}\raggedright
Depth MSE s1 / s2
\end{minipage} \\
\midrule\noalign{}
\endhead
\bottomrule\noalign{}
\endlastfoot
\texttt{paper\_maskoff} & \textbf{off} (uniform) & \textbf{0.734 /
0.842} & \textbf{0.650 / 0.773} & \textbf{0.230 / 0.352} &
\textbf{0.00259 / 0.00181} & 0.0426 / 0.0470 \\
\texttt{paper\_maskedge} & edge weighting only (\texttt{length\_scale}
0.02) & 0.611 / 0.655 & 0.485 / 0.541 & 0.184 / 0.206 & 0.00422 /
0.00358 & \textbf{0.0380 / 0.0383} \\
\texttt{paper\_maskblob} & blobs only (\texttt{center\_bias} 0) & 0.613
/ 0.673 & 0.487 / 0.571 & 0.180 / 0.197 & 0.00421 / 0.00332 & 0.0405 /
0.0394 \\
\texttt{paper\_maskboth} & blobs + edge (the old default) & 0.608 /
0.602 & 0.480 / 0.475 & 0.176 / 0.175 & 0.00435 / 0.00436 & 0.0390 /
0.0385 \\
\end{longtable}
}

Both \texttt{maskoff} seeds beat all six structured seeds on every PC
metric, and the gap (0.15--0.23 F1@0.03) is far larger than the seed
spread of any structured arm (0.005--0.06). Turning structured masking
off raises recall@0.03 from \textasciitilde0.48 to \textasciitilde0.71,
which is the same quantity the radial analysis below blamed on missing
leaf tips --- edge-biased masking spends the mask budget on the frame
border and starves exactly the outer plant. Blobs alone, edge bias alone
and both together are indistinguishable from each other.

\textbf{Robustness is the other half.} The same runs, scored under the
occluded validation pass --- one fixed reference corruption (procedural
leaves at 9--16 per image plus 1× sensor noise, eval seed 42), identical
for every run:

{\def\LTcaptype{none} % do not increment counter
\begin{longtable}[]{@{}
  >{\raggedright\arraybackslash}p{(\linewidth - 6\tabcolsep) * \real{0.2500}}
  >{\raggedright\arraybackslash}p{(\linewidth - 6\tabcolsep) * \real{0.2500}}
  >{\raggedright\arraybackslash}p{(\linewidth - 6\tabcolsep) * \real{0.2500}}
  >{\raggedright\arraybackslash}p{(\linewidth - 6\tabcolsep) * \real{0.2500}}@{}}
\toprule\noalign{}
\begin{minipage}[b]{\linewidth}\raggedright
arm
\end{minipage} & \begin{minipage}[b]{\linewidth}\raggedright
PC Chamfer clean s1 / s2
\end{minipage} & \begin{minipage}[b]{\linewidth}\raggedright
PC Chamfer occluded s1 / s2
\end{minipage} & \begin{minipage}[b]{\linewidth}\raggedright
degradation (occ / clean) s1 / s2
\end{minipage} \\
\midrule\noalign{}
\endhead
\bottomrule\noalign{}
\endlastfoot
\texttt{paper\_maskoff} & \textbf{0.00259 / 0.00181} & \textbf{0.00285 /
0.00199} & 1.10 / 1.10 \\
\texttt{paper\_maskedge} & 0.00422 / 0.00358 & 0.00432 / 0.00373 & 1.02
/ 1.04 \\
\texttt{paper\_maskblob} & 0.00421 / 0.00332 & 0.00433 / 0.00348 & 1.03
/ 1.05 \\
\texttt{paper\_maskboth} & 0.00435 / 0.00436 & 0.00444 / 0.00438 &
\textbf{1.02 / 1.00} \\
\texttt{paper\_paramsv2} & 0.00420 / 0.00435 & 0.00426 / 0.00435 &
\textbf{1.01 / 1.00} \\
\end{longtable}
}

Structured masking does what it was built for: corruption barely moves
its error (0--5 \%), while uniform masking loses 10 \%. At this
corruption strength, though, the robustness does not pay for itself ---
\texttt{maskoff} under corruption is still better than every structured
arm on \emph{clean} input. So the finding is a trade-off, not a win
either way, and it depends on how hard the test corruption is. Whether
the curves cross at stronger corruption is the open question
(\texttt{eval\_robustness.py}, below). The occluded pass logs Chamfer,
RGB and depth only --- F1/recall under corruption are not recorded yet.

\textbf{What structured masking already wins on, measured:}

{\def\LTcaptype{none} % do not increment counter
\begin{longtable}[]{@{}
  >{\raggedright\arraybackslash}p{(\linewidth - 4\tabcolsep) * \real{0.3333}}
  >{\raggedright\arraybackslash}p{(\linewidth - 4\tabcolsep) * \real{0.3333}}
  >{\raggedright\arraybackslash}p{(\linewidth - 4\tabcolsep) * \real{0.3333}}@{}}
\toprule\noalign{}
\begin{minipage}[b]{\linewidth}\raggedright
property
\end{minipage} & \begin{minipage}[b]{\linewidth}\raggedright
structured (10 runs: mask arms, loss arms, params v2)
\end{minipage} & \begin{minipage}[b]{\linewidth}\raggedright
\texttt{maskoff} (2 runs)
\end{minipage} \\
\midrule\noalign{}
\endhead
\bottomrule\noalign{}
\endlastfoot
PC Chamfer degradation under the reference corruption & 1.00--1.05× &
1.10× \\
Depth MSE (clean) & \textbf{0.0363--0.0407 --- every one of the 10 runs}
& 0.0426 / 0.0470 \\
F1@0.03 seed-to-seed spread & 0.006 (\texttt{maskboth}) & 0.108 \\
\end{longtable}
}

The depth result has no exceptions: all ten structured runs beat both
\texttt{maskoff} seeds. The stability gap is large but ambiguous --- a
model can be reproducible because it settles on a lower plateau. These
are supporting evidence; the headline claim needs one of the tests in
``Showing where structured masking helps'' below.

\textbf{Caveat before this is quoted:} \texttt{maskoff}'s two seeds are
0.734 vs 0.842, a much wider spread than any other arm, so the direction
is solid but the magnitude is not. A third \texttt{maskoff} seed is the
cheapest way to pin it (\textasciitilde10 h). Note also that depth goes
the other way: structured arms reach 0.038--0.040 Depth MSE against
0.043--0.047 for \texttt{maskoff}.

\paragraph{Spline params: v2 + pose/scale helps once its loss weight is
fixed}\label{spline-params-v2-posescale-helps-once-its-loss-weight-is-fixed}

Same masking (\texttt{maskboth}), same seeds; v1 is the
\texttt{paper\_maskboth} pair. \texttt{spline\_loss\_weight} is 5 for v1
and \textbf{0.6} for v2, which compensates for v2's z-scored targets
being \textasciitilde6× larger --- the confound that made the earlier
30k comparison read backwards.

{\def\LTcaptype{none} % do not increment counter
\begin{longtable}[]{@{}
  >{\raggedright\arraybackslash}p{(\linewidth - 8\tabcolsep) * \real{0.2000}}
  >{\raggedright\arraybackslash}p{(\linewidth - 8\tabcolsep) * \real{0.2000}}
  >{\raggedright\arraybackslash}p{(\linewidth - 8\tabcolsep) * \real{0.2000}}
  >{\raggedright\arraybackslash}p{(\linewidth - 8\tabcolsep) * \real{0.2000}}
  >{\raggedright\arraybackslash}p{(\linewidth - 8\tabcolsep) * \real{0.2000}}@{}}
\toprule\noalign{}
\begin{minipage}[b]{\linewidth}\raggedright
run
\end{minipage} & \begin{minipage}[b]{\linewidth}\raggedright
F1@0.03 s1 / s2
\end{minipage} & \begin{minipage}[b]{\linewidth}\raggedright
F1@0.01 s1 / s2
\end{minipage} & \begin{minipage}[b]{\linewidth}\raggedright
recall@0.03 s1 / s2
\end{minipage} & \begin{minipage}[b]{\linewidth}\raggedright
PC Chamfer s1 / s2
\end{minipage} \\
\midrule\noalign{}
\endhead
\bottomrule\noalign{}
\endlastfoot
v1 (\texttt{paper\_maskboth}) & 0.608 / 0.602 & 0.176 / 0.175 & 0.480 /
0.475 & 0.00435 / 0.00436 \\
\textbf{v2 + pose/scale} (\texttt{paper\_paramsv2}) & \textbf{0.619 /
0.615} & \textbf{0.224 / 0.215} & 0.483 / 0.483 & \textbf{0.00420 /
0.00435} \\
\end{longtable}
}

F1@0.01 improves by \textasciitilde25 \% (0.175 \(\to\) 0.219) against a
seed spread of 0.001--0.009, so this one is real. The gain is
concentrated at the tight threshold, i.e.~in placing points precisely,
which is what camera pose plus the normalisation radius should buy.

\paragraph{E7 --- point-cloud loss: QAL leads, Sinkhorn trades precision
for
recall}\label{e7-point-cloud-loss-qal-leads-sinkhorn-trades-precision-for-recall}

One seed each, masking fixed at \texttt{maskboth}; the QAL row is
\texttt{paper\_maskboth\_s1}.

{\def\LTcaptype{none} % do not increment counter
\begin{longtable}[]{@{}
  >{\raggedright\arraybackslash}p{(\linewidth - 12\tabcolsep) * \real{0.1429}}
  >{\raggedright\arraybackslash}p{(\linewidth - 12\tabcolsep) * \real{0.1429}}
  >{\raggedright\arraybackslash}p{(\linewidth - 12\tabcolsep) * \real{0.1429}}
  >{\raggedright\arraybackslash}p{(\linewidth - 12\tabcolsep) * \real{0.1429}}
  >{\raggedright\arraybackslash}p{(\linewidth - 12\tabcolsep) * \real{0.1429}}
  >{\raggedright\arraybackslash}p{(\linewidth - 12\tabcolsep) * \real{0.1429}}
  >{\raggedright\arraybackslash}p{(\linewidth - 12\tabcolsep) * \real{0.1429}}@{}}
\toprule\noalign{}
\begin{minipage}[b]{\linewidth}\raggedright
loss
\end{minipage} & \begin{minipage}[b]{\linewidth}\raggedright
F1@0.03
\end{minipage} & \begin{minipage}[b]{\linewidth}\raggedright
recall@0.03
\end{minipage} & \begin{minipage}[b]{\linewidth}\raggedright
precision@0.03
\end{minipage} & \begin{minipage}[b]{\linewidth}\raggedright
F1@0.01
\end{minipage} & \begin{minipage}[b]{\linewidth}\raggedright
PC Chamfer
\end{minipage} & \begin{minipage}[b]{\linewidth}\raggedright
Depth MSE
\end{minipage} \\
\midrule\noalign{}
\endhead
\bottomrule\noalign{}
\endlastfoot
QAL (t=0.01, \(\alpha\)=100) & \textbf{0.608} & 0.480 & 0.850 &
\textbf{0.176} & \textbf{0.00435} & 0.0390 \\
Chamfer (\texttt{pc\_loss\_weight} 10) & 0.578 & 0.450 & 0.846 & 0.144 &
0.00435 & \textbf{0.0363} \\
Sinkhorn only (blur 0.02, 2 048 pts) & 0.474 & \textbf{0.710} & 0.359 &
0.136 & 0.01032 & 0.0407 \\
\end{longtable}
}

Sinkhorn is the interesting one: it is the only loss that spreads
predictions over the whole plant (recall 0.71 vs 0.45--0.48) but it
places them badly (precision 0.36 vs 0.85). QAL and Chamfer do the
opposite. The combination (QAL with a Sinkhorn auxiliary term, which
\texttt{sinkhorn\_loss\_weight} already supports) is the obvious next
arm and has not been run.

\subsubsection{The 30 000-step sweep (screening
only)}\label{the-30-000-step-sweep-screening-only}

Everything below this line is the earlier 30 000-step pass. It sits
\emph{before} the F1 transition and its gaps were inside the seed gap,
so it is kept as screening evidence --- which settings to drop --- and
not as a result.

\textbf{Held fixed in every run --- including the edge weighting.} Token
masking: \texttt{model.structured\_mask} with \texttt{prob:\ 0.5} (half
the batches; the rest use the uniform shuffle),
\texttt{center\_bias:\ 1.5} (a smooth random field biased toward the
frame border, so masked blobs gather at the edges and the centre plant
stays visible) and one field shared across RGB, depth and the PC;
train-only, so validation always uses the uniform mask. Input occlusion:
procedural leaves entering from outside the frame
(\texttt{reach:\ {[}0.6,\ 1.2{]}}, 9--16 leaves, on every training
sample), interleaved in depth
(\texttt{depth\_quantile:\ {[}0.0,\ 0.8{]}}), plus sensor noise.
\textbf{No run in this screening pass varies \texttt{center\_bias},
\texttt{prob} or \texttt{reach}} --- that question is answered by the
600-epoch \texttt{paper\_mask*} arms above, and the answer is that
structured masking costs PC reconstruction.

\textbf{Against the three merge conditions above.} Split: same
\texttt{Sorghum\_15K} root yes. Global batch: 8 × 4 GPUs = 32 yes.
Budget in steps: 30 000 yes --- but that is a \emph{selection} schedule,
\textasciitilde15 \% of the programme's 197 400 (\texttt{e3\_3k}, same
94 steps/epoch, runs 2 100 epochs), so only the ranking transfers, not
the absolute numbers. The 3 000 plants are drawn by
\texttt{training.num\_train\_plants} with seed 42; that is \textbf{not}
the same draw as main's \texttt{data.max\_plants} /
\texttt{plant\_subset\_seed}, so it is not \texttt{e3\_3k}'s 3k.

\textbf{Noise floor --- read every table against this.} Training had no
seed until \texttt{training.seed} was added here, so \texttt{recipe\_b}
(unseeded) and \texttt{params\_v1\_seed1} (seed 1) are the \emph{same
config} run twice. Their gap is the only run-to-run error bar in this
section:

{\def\LTcaptype{none} % do not increment counter
\begin{longtable}[]{@{}
  >{\raggedright\arraybackslash}p{(\linewidth - 14\tabcolsep) * \real{0.1250}}
  >{\raggedright\arraybackslash}p{(\linewidth - 14\tabcolsep) * \real{0.1250}}
  >{\raggedright\arraybackslash}p{(\linewidth - 14\tabcolsep) * \real{0.1250}}
  >{\raggedright\arraybackslash}p{(\linewidth - 14\tabcolsep) * \real{0.1250}}
  >{\raggedright\arraybackslash}p{(\linewidth - 14\tabcolsep) * \real{0.1250}}
  >{\raggedright\arraybackslash}p{(\linewidth - 14\tabcolsep) * \real{0.1250}}
  >{\raggedright\arraybackslash}p{(\linewidth - 14\tabcolsep) * \real{0.1250}}
  >{\raggedright\arraybackslash}p{(\linewidth - 14\tabcolsep) * \real{0.1250}}@{}}
\toprule\noalign{}
\begin{minipage}[b]{\linewidth}\raggedright
\end{minipage} & \begin{minipage}[b]{\linewidth}\raggedright
val loss
\end{minipage} & \begin{minipage}[b]{\linewidth}\raggedright
val(occ) loss
\end{minipage} & \begin{minipage}[b]{\linewidth}\raggedright
PC Chamfer
\end{minipage} & \begin{minipage}[b]{\linewidth}\raggedright
F1@0.01
\end{minipage} & \begin{minipage}[b]{\linewidth}\raggedright
F1@0.03
\end{minipage} & \begin{minipage}[b]{\linewidth}\raggedright
EMD
\end{minipage} & \begin{minipage}[b]{\linewidth}\raggedright
Depth MSE
\end{minipage} \\
\midrule\noalign{}
\endhead
\bottomrule\noalign{}
\endlastfoot
seed-to-seed gap & 0.027 & 0.014 & 0.6 \% & 1.7 \% & 0.3 \% & 1.4 \% &
11 \% \\
\end{longtable}
}

\textbf{Which F1 threshold.} F1@0.03 is the headline column: its ceiling
is \textasciitilde1, so the value reads directly as a fraction of a
perfect reconstruction. F1@0.01 is kept beside it. Its ceiling is only
0.871 --- two independent 8 196-point draws of the \emph{same}
ground-truth cloud score that against each other (0.995 at 0.02, 0.999
at 0.03) --- so 0.154 is \textasciitilde18 \% of reachable, not 15 \%.
It is also slightly the more discriminating of the two here (recipe
spread / seed gap 3.1× vs 2.7×), and it is the column that exposes the
leaf-tip failure below. Do not move to a threshold where F1 \(\approx\)
0.9: every arm saturates there and the gaps compress.

\paragraph{E6 --- structured-mask blob size × sensor
noise}\label{e6-structured-mask-blob-size-sensor-noise}

{\def\LTcaptype{none} % do not increment counter
\begin{longtable}[]{@{}
  >{\raggedright\arraybackslash}p{(\linewidth - 18\tabcolsep) * \real{0.1000}}
  >{\raggedright\arraybackslash}p{(\linewidth - 18\tabcolsep) * \real{0.1000}}
  >{\raggedright\arraybackslash}p{(\linewidth - 18\tabcolsep) * \real{0.1000}}
  >{\raggedright\arraybackslash}p{(\linewidth - 18\tabcolsep) * \real{0.1000}}
  >{\raggedright\arraybackslash}p{(\linewidth - 18\tabcolsep) * \real{0.1000}}
  >{\raggedright\arraybackslash}p{(\linewidth - 18\tabcolsep) * \real{0.1000}}
  >{\raggedright\arraybackslash}p{(\linewidth - 18\tabcolsep) * \real{0.1000}}
  >{\raggedright\arraybackslash}p{(\linewidth - 18\tabcolsep) * \real{0.1000}}
  >{\raggedright\arraybackslash}p{(\linewidth - 18\tabcolsep) * \real{0.1000}}
  >{\raggedright\arraybackslash}p{(\linewidth - 18\tabcolsep) * \real{0.1000}}@{}}
\toprule\noalign{}
\begin{minipage}[b]{\linewidth}\raggedright
run
\end{minipage} & \begin{minipage}[b]{\linewidth}\raggedright
blob size (\texttt{length\_scale})
\end{minipage} & \begin{minipage}[b]{\linewidth}\raggedright
train noise (rgb/depth/pc)
\end{minipage} & \begin{minipage}[b]{\linewidth}\raggedright
val loss
\end{minipage} & \begin{minipage}[b]{\linewidth}\raggedright
val(occ) loss
\end{minipage} & \begin{minipage}[b]{\linewidth}\raggedright
PC Chamfer
\end{minipage} & \begin{minipage}[b]{\linewidth}\raggedright
F1@0.03
\end{minipage} & \begin{minipage}[b]{\linewidth}\raggedright
F1@0.01
\end{minipage} & \begin{minipage}[b]{\linewidth}\raggedright
RGB MSE
\end{minipage} & \begin{minipage}[b]{\linewidth}\raggedright
Depth MSE
\end{minipage} \\
\midrule\noalign{}
\endhead
\bottomrule\noalign{}
\endlastfoot
\texttt{recipe\_a\_ls035\_noise1x} & 0.35 & 1× (0.03 / 0.015 / 0.008) &
0.5829 & 0.5964 & 0.004703 & 0.5832 & 0.1582 & 0.6738 & 0.03350 \\
\texttt{recipe\_b\_ls020\_noise1x} & 0.20 & 1× & 0.5623 & 0.5853 &
0.004644 & 0.5856 & 0.1567 & 0.6649 & 0.03355 \\
\texttt{recipe\_c\_ls020\_noise0p5x} & 0.20 & 0.5× & 0.5908 & 0.5993 &
0.004716 & 0.5816 & 0.1536 & 0.6715 & 0.03219 \\
\texttt{recipe\_d\_ls020\_noise2x} & 0.20 & 2× & 0.5891 & 0.6010 &
0.004680 & 0.5817 & 0.1503 & 0.6610 & 0.03221 \\
\texttt{params\_v1\_seed1} (= B, seed 1) & 0.20 & 1× & 0.5891 & 0.5996 &
0.004670 & 0.5841 & 0.1541 & 0.6773 & 0.02997 \\
\end{longtable}
}

Nothing is bold: no gap in this table exceeds the seed-to-seed gap.

B has the best clean and occluded val loss, but \textbf{B run a second
time lands behind A and level with C and D}: every recipe gap (A-B
0.021, C-B 0.029, D-B 0.027 val loss) is the size of the seed gap. The
sweep does not rank these recipes --- and the order of A and B flips
with the F1 threshold (A \textgreater{} B at 0.01, B \textgreater{} A at
0.02 and 0.03). The default stays B (\texttt{length\_scale\ 0.20}, 1×
noise) --- best observed, not shown better. Separating them needs
\(\ge\) 3 seeds per arm.

\textbf{Where the model fails: leaf tips, not the frame edge.} Trained
recipe-B model, 12 val plants, 4 mask draws; radius 0 = image centre, 1
= corner (for the PC, the xy radius of the unit-normalised cloud).
Errors on \emph{masked} patches:

{\def\LTcaptype{none} % do not increment counter
\begin{longtable}[]{@{}
  >{\raggedright\arraybackslash}p{(\linewidth - 10\tabcolsep) * \real{0.1667}}
  >{\raggedright\arraybackslash}p{(\linewidth - 10\tabcolsep) * \real{0.1667}}
  >{\raggedright\arraybackslash}p{(\linewidth - 10\tabcolsep) * \real{0.1667}}
  >{\raggedright\arraybackslash}p{(\linewidth - 10\tabcolsep) * \real{0.1667}}
  >{\raggedright\arraybackslash}p{(\linewidth - 10\tabcolsep) * \real{0.1667}}
  >{\raggedright\arraybackslash}p{(\linewidth - 10\tabcolsep) * \real{0.1667}}@{}}
\toprule\noalign{}
\begin{minipage}[b]{\linewidth}\raggedright
radius
\end{minipage} & \begin{minipage}[b]{\linewidth}\raggedright
0--0.2
\end{minipage} & \begin{minipage}[b]{\linewidth}\raggedright
0.2--0.4
\end{minipage} & \begin{minipage}[b]{\linewidth}\raggedright
0.4--0.6
\end{minipage} & \begin{minipage}[b]{\linewidth}\raggedright
0.6--0.8
\end{minipage} & \begin{minipage}[b]{\linewidth}\raggedright
0.8--1
\end{minipage} \\
\midrule\noalign{}
\endhead
\bottomrule\noalign{}
\endlastfoot
plant share of patches & 71 \% & 40 \% & 16 \% & 3 \% & 0.2 \% \\
P(masked), uniform mask & 0.64 & 0.61 & 0.60 & 0.60 & 0.62 \\
P(masked), structured (\texttt{center\_bias} 1.5) & 0.64 & 0.70 & 0.81 &
0.86 & 0.89 \\
RGB MSE, all masked patches & 0.18 & 0.17 & 0.09 & 0.02 & 0.002 \\
RGB MSE, plant patches only & 0.20 & 0.25 & 0.29 & 0.27 & 0.60 \\
PC: target point \(\to\) nearest prediction & 0.025 & 0.042 & 0.072 &
0.119 & 0.233 \\
\end{longtable}
}

The edge bias does what it says --- masking rises from 0.64 to 0.89
toward the corner --- but the corner is background, which is why the
error there is near zero. On the plant itself error \emph{grows}
outward, and PC error on the outermost points is \textasciitilde10× the
centre's: the model misses leaf tips. That is the low F1@0.01 (precision
0.40, recall 0.11 at the end of training --- predicted points sit on the
dense core and do not reach the tips). So much of the edge-biased mask
budget is spent on background; masking the plant's \emph{outer parts}
would target the actual failure. The 600-epoch arms above confirm the
mechanism from the other side: switching structured masking off raises
recall@0.03 from \textasciitilde0.48 to \textasciitilde0.71. Preview of
the masks: \texttt{reports/mask\_preview\_recipe\_b.png},
\texttt{reports/mask\_blob\_size.png} (uniform vs blob sizes at a fixed
mask count) and \texttt{reports/noise\_levels.png} (the 0.5× / 1× / 2×
sensor noise).

\paragraph{Spline params: do they help, and does fixing their encoding
help
more?}\label{spline-params-do-they-help-and-does-fixing-their-encoding-help-more}

Measured on the data before changing anything: 5 of the 9 plant dims
(\texttt{panicle\_*}) are constant; the per-sample PC normalisation
removes absolute size (corr(\texttt{stem\_length}, cloud height) is 0.68
raw, 0.08 after normalisation); and the params are in the plant's world
frame while the target is in the camera frame (\texttt{*\_nc\_cam.ply}
is exactly \texttt{worldToCamera\ @\ world}), with the
\texttt{camera\_pose.json} that links them never read.
\texttt{model.param\_encoding:\ v2} (z-scored, sin/cos angles, constants
dropped, leaf count added) and \texttt{model.geometry\_cond:\ true}
(camera rotation + normalisation radius as an always-visible token)
address those. Same seed, only the params path differs:

{\def\LTcaptype{none} % do not increment counter
\begin{longtable}[]{@{}
  >{\raggedright\arraybackslash}p{(\linewidth - 14\tabcolsep) * \real{0.1250}}
  >{\raggedright\arraybackslash}p{(\linewidth - 14\tabcolsep) * \real{0.1250}}
  >{\raggedright\arraybackslash}p{(\linewidth - 14\tabcolsep) * \real{0.1250}}
  >{\raggedright\arraybackslash}p{(\linewidth - 14\tabcolsep) * \real{0.1250}}
  >{\raggedright\arraybackslash}p{(\linewidth - 14\tabcolsep) * \real{0.1250}}
  >{\raggedright\arraybackslash}p{(\linewidth - 14\tabcolsep) * \real{0.1250}}
  >{\raggedright\arraybackslash}p{(\linewidth - 14\tabcolsep) * \real{0.1250}}
  >{\raggedright\arraybackslash}p{(\linewidth - 14\tabcolsep) * \real{0.1250}}@{}}
\toprule\noalign{}
\begin{minipage}[b]{\linewidth}\raggedright
run
\end{minipage} & \begin{minipage}[b]{\linewidth}\raggedright
PC Chamfer
\end{minipage} & \begin{minipage}[b]{\linewidth}\raggedright
F1@0.01
\end{minipage} & \begin{minipage}[b]{\linewidth}\raggedright
F1@0.02
\end{minipage} & \begin{minipage}[b]{\linewidth}\raggedright
F1@0.03
\end{minipage} & \begin{minipage}[b]{\linewidth}\raggedright
EMD
\end{minipage} & \begin{minipage}[b]{\linewidth}\raggedright
RGB MSE
\end{minipage} & \begin{minipage}[b]{\linewidth}\raggedright
Depth MSE
\end{minipage} \\
\midrule\noalign{}
\endhead
\bottomrule\noalign{}
\endlastfoot
\texttt{params\_v1\_seed1} & \textbf{0.004670} & 0.1541 & 0.4097 &
\textbf{0.5841} & 0.2966 & 0.6773 & 0.02997 \\
\texttt{params\_v2cond\_seed1} & 0.004967 & \textbf{0.1587} & 0.4057 &
0.5748 & 0.2948 & 0.6758 & 0.03256 \\
\end{longtable}
}

Bold = better by more than the seed-to-seed gap; every other gap is
within noise.

\texttt{eval\_param\_oracle.py} then asks each model directly: all param
tokens visible, the sample's own params vs the val-set mean; identical
RGB/depth/PC masks in both passes. Positive = worse without the sample's
own params:

{\def\LTcaptype{none} % do not increment counter
\begin{longtable}[]{@{}
  >{\raggedright\arraybackslash}p{(\linewidth - 12\tabcolsep) * \real{0.1429}}
  >{\raggedright\arraybackslash}p{(\linewidth - 12\tabcolsep) * \real{0.1429}}
  >{\raggedright\arraybackslash}p{(\linewidth - 12\tabcolsep) * \real{0.1429}}
  >{\raggedright\arraybackslash}p{(\linewidth - 12\tabcolsep) * \real{0.1429}}
  >{\raggedright\arraybackslash}p{(\linewidth - 12\tabcolsep) * \real{0.1429}}
  >{\raggedright\arraybackslash}p{(\linewidth - 12\tabcolsep) * \real{0.1429}}
  >{\raggedright\arraybackslash}p{(\linewidth - 12\tabcolsep) * \real{0.1429}}@{}}
\toprule\noalign{}
\begin{minipage}[b]{\linewidth}\raggedright
checkpoint
\end{minipage} & \begin{minipage}[b]{\linewidth}\raggedright
PC Chamfer
\end{minipage} & \begin{minipage}[b]{\linewidth}\raggedright
F1@0.01
\end{minipage} & \begin{minipage}[b]{\linewidth}\raggedright
F1@0.02
\end{minipage} & \begin{minipage}[b]{\linewidth}\raggedright
RGB MSE
\end{minipage} & \begin{minipage}[b]{\linewidth}\raggedright
Depth MSE
\end{minipage} & \begin{minipage}[b]{\linewidth}\raggedright
occluded PC Chamfer
\end{minipage} \\
\midrule\noalign{}
\endhead
\bottomrule\noalign{}
\endlastfoot
v1 & +7.0 \% & +1.4 \% & +2.1 \% & -1.7 \% & +7.1 \% & +7.9 \% \\
v2 + pose/scale & +4.6 \% & +5.3 \% & +2.5 \% & +0.4 \% & -0.2 \% & +4.5
\% \\
\end{longtable}
}

\textbf{The params already help in v1} (7 \% Chamfer, well above the 0.6
\% noise floor). v2 + pose/scale looked 6 \% \emph{worse} here, and the
suspected cause --- v2's z-scored Smooth-L1 targets being
\textasciitilde6× larger at the same \texttt{spline\_loss\_weight:\ 5},
so the param term steals gradient from the PC term --- turned out to be
right: \textbf{at \texttt{spline\_loss\_weight:\ 0.6} and 600 epochs v2
wins} (F1@0.01 0.219 vs 0.175), see the paper-run section above. The
oracle numbers below still stand; they are measured within a single
checkpoint.

\paragraph{E7 --- point-cloud loss study
(running)}\label{e7-point-cloud-loss-study-running}

Loss magnitudes measured on the trained recipe-B model set the weights:
QAL 0.060, Chamfer 0.0053 (11× smaller), Sinkhorn \textasciitilde0.055.
Chamfer therefore needs \texttt{pc\_loss\_weight} \(\approx\) 10 to pull
as hard as QAL at 1. \texttt{loss\_name:\ sinkhorn} means Sinkhorn
\emph{alone} (Sinkhorn was only an auxiliary term before), on 2 048
sampled points.

{\def\LTcaptype{none} % do not increment counter
\begin{longtable}[]{@{}
  >{\raggedright\arraybackslash}p{(\linewidth - 4\tabcolsep) * \real{0.3333}}
  >{\raggedright\arraybackslash}p{(\linewidth - 4\tabcolsep) * \real{0.3333}}
  >{\raggedright\arraybackslash}p{(\linewidth - 4\tabcolsep) * \real{0.3333}}@{}}
\toprule\noalign{}
\begin{minipage}[b]{\linewidth}\raggedright
loss
\end{minipage} & \begin{minipage}[b]{\linewidth}\raggedright
arms
\end{minipage} & \begin{minipage}[b]{\linewidth}\raggedright
runs
\end{minipage} \\
\midrule\noalign{}
\endhead
\bottomrule\noalign{}
\endlastfoot
QAL & t = 0.005 / \textbf{0.01} / 0.02 at \(\alpha\) = 100; \(\alpha\) =
30 / 300 at t = 0.01 & \texttt{loss\_q\_t005},
\texttt{params\_v1\_seed1} (t=0.01, \(\alpha\)=100),
\texttt{loss\_q\_t020}, \texttt{loss\_q\_a030},
\texttt{loss\_q\_a300} \\
Chamfer & \texttt{pc\_loss\_weight} 3 / 10 / 30 & \texttt{loss\_c\_w03},
\texttt{loss\_c\_w10}, \texttt{loss\_c\_w30} \\
Sinkhorn only & blur 0.01 / 0.02 / 0.05 at weight 1; blur 0.02 at weight
3 & \texttt{loss\_s\_b010\_w1}, \texttt{loss\_s\_b020\_w1},
\texttt{loss\_s\_b050\_w1}, \texttt{loss\_s\_b020\_w3} \\
\end{longtable}
}

Job array 15884575 (\texttt{train\_4m\_loss.sbatch}); the Sinkhorn-only
path passed its smoke test first. The array allows two arms at a time,
but with the mech-ai account at its 17-GPU cap it has run one at a time
(\textasciitilde5 h per arm). Read it with
\texttt{python\ compare\_recipes.py\ -\/-loss-study}, which shows only
metrics computed identically for every arm --- val loss is a different
quantity in each.

Finished so far (2026-09-22): 1 of 11 new arms.

{\def\LTcaptype{none} % do not increment counter
\begin{longtable}[]{@{}
  >{\raggedright\arraybackslash}p{(\linewidth - 18\tabcolsep) * \real{0.1000}}
  >{\raggedright\arraybackslash}p{(\linewidth - 18\tabcolsep) * \real{0.1000}}
  >{\raggedright\arraybackslash}p{(\linewidth - 18\tabcolsep) * \real{0.1000}}
  >{\raggedright\arraybackslash}p{(\linewidth - 18\tabcolsep) * \real{0.1000}}
  >{\raggedright\arraybackslash}p{(\linewidth - 18\tabcolsep) * \real{0.1000}}
  >{\raggedright\arraybackslash}p{(\linewidth - 18\tabcolsep) * \real{0.1000}}
  >{\raggedright\arraybackslash}p{(\linewidth - 18\tabcolsep) * \real{0.1000}}
  >{\raggedright\arraybackslash}p{(\linewidth - 18\tabcolsep) * \real{0.1000}}
  >{\raggedright\arraybackslash}p{(\linewidth - 18\tabcolsep) * \real{0.1000}}
  >{\raggedright\arraybackslash}p{(\linewidth - 18\tabcolsep) * \real{0.1000}}@{}}
\toprule\noalign{}
\begin{minipage}[b]{\linewidth}\raggedright
run
\end{minipage} & \begin{minipage}[b]{\linewidth}\raggedright
loss
\end{minipage} & \begin{minipage}[b]{\linewidth}\raggedright
setting
\end{minipage} & \begin{minipage}[b]{\linewidth}\raggedright
PC Chamfer
\end{minipage} & \begin{minipage}[b]{\linewidth}\raggedright
F1@0.01
\end{minipage} & \begin{minipage}[b]{\linewidth}\raggedright
F1@0.02
\end{minipage} & \begin{minipage}[b]{\linewidth}\raggedright
F1@0.03
\end{minipage} & \begin{minipage}[b]{\linewidth}\raggedright
EMD
\end{minipage} & \begin{minipage}[b]{\linewidth}\raggedright
RGB MSE
\end{minipage} & \begin{minipage}[b]{\linewidth}\raggedright
Depth MSE
\end{minipage} \\
\midrule\noalign{}
\endhead
\bottomrule\noalign{}
\endlastfoot
\texttt{params\_v1\_seed1} & QAL & t=0.01, \(\alpha\)=100 & 0.004670 &
0.1541 & 0.4097 & 0.5841 & 0.2966 & 0.6773 & 0.02997 \\
\texttt{loss\_q\_t005} & QAL & t=0.005, \(\alpha\)=100 & 0.004700 &
0.1540 & 0.4060 & 0.5822 & 0.2970 & 0.6910 & 0.02898 \\
\end{longtable}
}

Halving the QAL threshold moves nothing past the noise floor.

\paragraph{What this branch adds, and what will conflict with
main}\label{what-this-branch-adds-and-what-will-conflict-with-main}

\begin{itemize}
\tightlist
\item
  Config keys: \texttt{training.num\_train\_plants},
  \texttt{training.views\_per\_epoch}, \texttt{training.max\_steps},
  \texttt{training.seed}, \texttt{occlusion.eval\_occlusion},
  \texttt{model.param\_encoding}, \texttt{model.geometry\_cond}, and
  \texttt{loss\_name:\ sinkhorn}. All default off; with none set, a
  model is bit-identical to before (same parameters, init and loss).
\item
  Main implements plant subsetting and view sampling separately
  (\texttt{data.max\_plants}, \texttt{data.view\_sampling}). Keep one
  when merging; they will conflict in \texttt{train\_sorghum\_4m.py} and
  the datasets.
\item
  New files, at the root because this branch predates main's
  reorganisation: \texttt{view\_sampler.py},
  \texttt{compare\_recipes.py}, \texttt{eval\_param\_oracle.py},
  \texttt{compute\_param\_stats.py},
  \texttt{config\_4m\_\{recipe,params,loss\}\_*.yaml},
  \texttt{train\_4m\_\{recipe,params,loss\}.sbatch}. On merge they
  belong in \texttt{configs/}, \texttt{slurm/} and \texttt{eval/}.
\item
  Slurm: submitted under \texttt{-\/-account=mech-ai}, whose GPU cap
  (17) is shared across the lab. Arrays run \textbf{two at a time = 8
  GPUs}, the ceiling this project keeps.
\item
  Launchers: \texttt{train\_4m\_paper.sbatch} (the twelve 600-epoch
  runs), \texttt{train\_4m\_recipe.sbatch} /
  \texttt{train\_4m\_loss.sbatch} (the 30k screening),
  \texttt{train\_4m\_continue.sbatch} (resume an existing run at a
  larger \texttt{-\/-max\_steps}; unused now that the runs are
  single-cosine).
\end{itemize}

\paragraph{Showing where structured masking helps --- the three
legitimate
tests}\label{showing-where-structured-masking-helps-the-three-legitimate-tests}

The claim the method needs is ``structured masking makes the model
better at the thing it is for''. Clean reconstruction is not that thing,
and on it the answer is no. Three tests could show a real advantage;
each is worth running, and if none shows one, the honest write-up is the
trade-off above.

\begin{enumerate}
\def\labelenumi{\arabic{enumi}.}
\tightlist
\item
  \textbf{Stronger corruption and test-time blob masks, on TEST} ---
  \texttt{eval\_robustness.py\ -\/-split\ test} scores
  \texttt{best\_model.pth} (chosen on val loss, so val numbers carry
  selection bias; report test) on one view per test plant (2 250 plants,
  view fixed by seed --- 10× the 225-view val subset, no plant counted
  twice). Levels: \texttt{clean}, \texttt{1x}, \texttt{noise2x},
  \texttt{noise4x}, \texttt{leaves2x}, \texttt{leaves2x\_noise2x}, and
  two \textbf{blob-mask} levels, \texttt{blobmask} (clean input) and
  \texttt{blobmask\_1x} (1× corruption), which impose the
  \texttt{maskboth} blob geometry on \emph{every} model's token mask at
  test time. Every other evaluation masks tokens uniformly --- the
  regime \texttt{maskoff} trained on --- so the blob levels are the
  first test on structured masking's own terms: filling large contiguous
  holes, the shape real occlusion takes. Queued as two 1-GPU jobs
  (\texttt{.smoke/eval\_robustness.sbatch}, \texttt{EVAL\_RUNS} /
  \texttt{EVAL\_OUT}, results in
  \texttt{reports/robustness\_test\_\{A,B\}.json}), after the 1
  000-epoch continuation.
\item
  \textbf{The programme's actual metric} --- decision 6.4 scores
  value-add with a linear probe on height, leaf angle, leaf count and
  biomass, \emph{not} reconstruction. Structured masking could help the
  representation even while it hurts reconstruction; the probe is the
  test (see ``Two gaps'').
\item
  \textbf{Real data (E10)} --- real occlusion is what the synthetic
  corruption imitates; an advantage there is the strongest possible
  evidence.
\end{enumerate}

\paragraph{Operational notes
(2026-09-28)}\label{operational-notes-2026-09-28}

\begin{itemize}
\tightlist
\item
  \textbf{The \texttt{/work/mech-ai-scratch} filesystem is degraded,
  cluster-wide.} It is an NFS mount
  (\texttt{novastor010:/stor010/}\allowbreak\texttt{mech-ai-scratch}) at
  \textbf{99 \% full --- 150 of 152 TB}. On a compute node
  (\texttt{nova18-wide-10}), \texttt{import\ numpy} took 175 s and
  \texttt{import\ torch}, \texttt{import\ open3d} and listing 500
  dataset folders each ran past 300 s; \texttt{ls} takes 4.6 s in
  \texttt{/work/mech-ai-scratch/yongyun} and 12 s in the dataset root,
  against 4 ms in \texttt{/home} and 0.7 s in \texttt{/work/mech-ai}.
  This repo's own checkpoints are \textasciitilde0.6 TB
  (\textasciitilde0.4 \% of the used space), so freeing them is
  housekeeping, not a fix; the fix is HPC support and the lab freeing
  space. Symptoms so far: continuation task 16613000\_0 held 4 GPUs for
  7 h without leaving the dataset scan; the first robustness eval hit
  its 4 h limit the same way. \textbf{Do not submit GPU jobs until
  \texttt{time\ python\ -c\ "import\ torch"} is back under a minute} ---
  a stalled job holds 4 GPUs doing nothing.
\item
  \textbf{Auto-submit watcher:} \texttt{.smoke/watch\_and\_submit.sh}
  (log in \texttt{.smoke/watch.log}) probes \texttt{import\ torch} every
  10 min; after two healthy probes it pre-builds the index cache with
  \texttt{.smoke/build\_index\_cache.py} (torch-free, same keys as the
  loaders, val/test/train), then submits the twelve-run continuation to
  1 000 epochs (8 GPUs) and the two test-split robustness evals after
  it.
\item
  \textbf{Dataset folder index is now cached} in \texttt{.index\_cache/}
  (per split, per 3M/4M). The first run pays one scan; every later run
  and evaluation reads the cache. \texttt{SORGHUM\_INDEX\_REFRESH=1}
  forces a rescan.
\item
  \textbf{\texttt{train\_4m\_continue.sbatch} takes \texttt{RUNS}
  colon-separated.} \texttt{sbatch\ -\/-export} splits its argument on
  commas, so a comma list silently delivered only the first run and
  array job 16613000 lost eleven of twelve tasks.
\item
  \textbf{Continuation checkpoints every 100 epochs}
  (\texttt{SAVE\_FREQ}, default 100 in
  \texttt{train\_4m\_continue.sbatch}): at 1.3 GB each, the configs'
  30-epoch cadence would add \textasciitilde200 GB to a 99 \%-full disk;
  100 adds \textasciitilde60 GB.
\item
  \textbf{Pending:} all twelve \texttt{paper\_*} runs resumed from epoch
  600 to 1 000 (\texttt{TARGET\_STEPS=94000}), then
  \texttt{eval\_robustness.py\ -\/-split\ test} on the results.
\end{itemize}

\paragraph{Checkpoint retention --- what the paper
needs}\label{checkpoint-retention-what-the-paper-needs}

Keep, per run: \textbf{\texttt{best\_model.pth}} (the reported numbers
come from it --- it is not necessarily the last epoch), \textbf{the
latest checkpoint} until the run is finished (it is the resume point),
and \textbf{\texttt{training\_history.json} + \texttt{config.json}}
always (kilobytes; they back every curve and every ``what did this run
use'' question). Intermediate checkpoints are not needed --- the val
curves in the history replace them. By group:

{\def\LTcaptype{none} % do not increment counter
\begin{longtable}[]{@{}
  >{\raggedright\arraybackslash}p{(\linewidth - 4\tabcolsep) * \real{0.3333}}
  >{\raggedright\arraybackslash}p{(\linewidth - 4\tabcolsep) * \real{0.3333}}
  >{\raggedright\arraybackslash}p{(\linewidth - 4\tabcolsep) * \real{0.3333}}@{}}
\toprule\noalign{}
\begin{minipage}[b]{\linewidth}\raggedright
runs
\end{minipage} & \begin{minipage}[b]{\linewidth}\raggedright
keep
\end{minipage} & \begin{minipage}[b]{\linewidth}\raggedright
delete when space is needed
\end{minipage} \\
\midrule\noalign{}
\endhead
\bottomrule\noalign{}
\endlastfoot
\texttt{paper\_*} (the results) & best + latest (epoch 600 now, 1 000
after the continuation) + history & intermediates (19 per run,
\textasciitilde297 GB) \\
\texttt{params\_v1\_seed1}, \texttt{params\_v2cond\_seed1} & best, if
the oracle table is used & intermediates \\
30k screening (\texttt{recipe\_*}, \texttt{loss\_*}) & history & all
weights (\textasciitilde157 GB) \\
\texttt{\_aborted\_*} & history (supports ``600 epochs is enough'') &
all weights (\textasciitilde14 GB) \\
\end{longtable}
}

Two cautions: other project folders under
\texttt{/work/mech-ai-scratch/yongyun} (e.g.~\texttt{embodiedmae4m},
linked to the epoch-760 results) may hold models the paper cites ---
check before deleting there. And scratch is not backed up, so copy the
handful of \texttt{best\_model.pth} files the paper reports off scratch
(\texttt{/work/mech-ai} or external storage). Nothing has been deleted
yet.

\paragraph{What to run next, in order}\label{what-to-run-next-in-order}

\begin{enumerate}
\def\labelenumi{\arabic{enumi}.}
\tightlist
\item
  \textbf{A third \texttt{maskoff} seed} --- its two seeds are 0.734 and
  0.842 on F1@0.03, the widest spread of any arm, and it carries the
  headline claim (\textasciitilde10 h).
\item
  \textbf{QAL + Sinkhorn auxiliary} (\texttt{sinkhorn\_loss\_weight}
  \textgreater{} 0 on top of \texttt{loss\_name:\ qal\_loss}): Sinkhorn
  alone reaches recall 0.71 at precision 0.36, QAL alone 0.48 at 0.85,
  so the combination is the obvious arm.
\item
  \textbf{v2 params + \texttt{maskoff}} --- both wins are measured
  separately, never together.
\end{enumerate}

\subsection{Environment}\label{environment}

Conda env name is \texttt{det} (Python 3.12, PyTorch 2.5 + CUDA 12.4;
see \texttt{environment.yml}). On Nova it lives at
\texttt{/work/mech-ai/alloy/.conda/envs/det}. Activate with
\texttt{conda\ activate\ det} before running anything.

On Blackwell / sm\_120 GPUs (RTX PRO 6000) \texttt{det} will not work
--- those need the CUDA 12.8 build in the separate \texttt{det\_cu128}
env.

\subsection{Common commands}\label{common-commands}

All commands assume \texttt{cwd\ =\ repo\ root} and the env active.

\subsubsection{Train 4M (the main entry
point)}\label{train-4m-the-main-entry-point}

\begin{Shaded}
\begin{Highlighting}[]
\CommentTok{\# Single GPU}
\ExtensionTok{python}\NormalTok{ train\_sorghum\_4m.py }\AttributeTok{{-}{-}config}\NormalTok{ configs/config\_4m.yaml }\AttributeTok{{-}{-}world\_size}\NormalTok{ 1}

\CommentTok{\# Multi{-}GPU via torchrun (preferred — sets LOCAL\_RANK/RANK/WORLD\_SIZE)}
\ExtensionTok{torchrun} \AttributeTok{{-}{-}standalone} \AttributeTok{{-}{-}nproc\_per\_node}\OperatorTok{=}\NormalTok{4 train\_sorghum\_4m.py }\AttributeTok{{-}{-}config}\NormalTok{ configs/config\_4m.yaml}

\CommentTok{\# Multi{-}GPU via mp.spawn fallback (set distributed.world\_size in the YAML)}
\ExtensionTok{python}\NormalTok{ train\_sorghum\_4m.py }\AttributeTok{{-}{-}config}\NormalTok{ configs/config\_4m.yaml}
\end{Highlighting}
\end{Shaded}

\subsubsection{Train 3M}\label{train-3m}

\begin{Shaded}
\begin{Highlighting}[]
\ExtensionTok{python}\NormalTok{ train\_sorghum.py }\AttributeTok{{-}{-}config}\NormalTok{ configs/config.yaml                  }\CommentTok{\# single GPU}
\ExtensionTok{python}\NormalTok{ train\_sorghum\_multi.py }\AttributeTok{{-}{-}config}\NormalTok{ configs/config.yaml            }\CommentTok{\# mp.spawn}
\end{Highlighting}
\end{Shaded}

\subsubsection{Cross-modal distillation}\label{cross-modal-distillation}

\texttt{train\_sorghum\_4m\_distill.py} trains the 4M model so that
\textbf{any single modality, with the other three fully masked,
reconstructs all four} --- a frozen full-modal teacher supplies
token-aligned decoder features and a CLS latent as the privileged
target, and the student is warm-started from the same checkpoint.

\begin{Shaded}
\begin{Highlighting}[]
\ExtensionTok{python}\NormalTok{ train\_sorghum\_4m\_distill.py }\AttributeTok{{-}{-}config}\NormalTok{ configs/config\_4m\_distill\_15k\_rgb2pc.yaml}
\end{Highlighting}
\end{Shaded}

\subsubsection{Validate / dump
reconstructions}\label{validate-dump-reconstructions}

\begin{Shaded}
\begin{Highlighting}[]
\ExtensionTok{python}\NormalTok{ validate.py }\AttributeTok{{-}{-}checkpoint}\NormalTok{ outputs/}\OperatorTok{\textless{}}\NormalTok{run}\OperatorTok{\textgreater{}}\NormalTok{/best\_model.pth }\DataTypeTok{\textbackslash{}}
                   \AttributeTok{{-}{-}output\_dir}\NormalTok{ vis\_val }\AttributeTok{{-}{-}config}\NormalTok{ configs/config.yaml }\AttributeTok{{-}{-}num\_samples}\NormalTok{ 6}
\end{Highlighting}
\end{Shaded}

\subsubsection{Sanity-check a dataset
folder}\label{sanity-check-a-dataset-folder}

\begin{Shaded}
\begin{Highlighting}[]
\ExtensionTok{python}\NormalTok{ sorghum\_dataset.py    /path/to/Sorghum\_15K    }\CommentTok{\# 3M}
\ExtensionTok{python}\NormalTok{ sorghum\_dataset\_4m.py /path/to/Sorghum\_15K    }\CommentTok{\# 4M (requires *\_spline.yml)}
\end{Highlighting}
\end{Shaded}

\subsubsection{Smoke-test the model
definitions}\label{smoke-test-the-model-definitions}

\begin{Shaded}
\begin{Highlighting}[]
\ExtensionTok{python}\NormalTok{ embodied\_mae.py       }\CommentTok{\# builds embodied\_mae\_base, one forward pass on dummy tensors}
\ExtensionTok{python}\NormalTok{ embodied\_mae\_4m.py    }\CommentTok{\# same for 4M}
\end{Highlighting}
\end{Shaded}

There is no test suite, no linter, and no Makefile.

\subsection{Configuration model}\label{configuration-model}

Both training entry points layer config in this order: YAML \(\to\) CLI
flags \(\to\) defaults. The YAML is the source of truth; CLI flags only
override specific keys. Notable keys:

\begin{itemize}
\tightlist
\item
  \texttt{data.data\_root} expects
  \texttt{\textless{}root\textgreater{}/train/},
  \texttt{\textless{}root\textgreater{}/val/} and
  \texttt{\textless{}root\textgreater{}/test/} siblings, each containing
  one folder per sample (see ``Dataset layout'').
\item
  \texttt{data.view\_sampling:\ true} --- each plant contributes
  \textbf{one randomly chosen view per epoch}, so an epoch over 105 000
  train samples is 10 500 items, not 105 000. \texttt{view\_seed} must
  be identical across arms of an ablation so they see the same views in
  the same order.
\item
  \texttt{data.max\_plants} --- caps the \textbf{train} split at a
  nested random subset of N plants (\texttt{plant\_subset\_seed} fixes
  the draw); \texttt{null} uses all of them. This is what E3 varies. The
  subsets nest, so 1k \(\subset\) 3k \(\subset\) 10k and a step of the
  curve is never partly a change of sample composition. It selects
  \textbf{plants, keeping all ten views} --- dropping views would change
  the augmentation regime that decision 6.1 fixes, which is E5's
  question, not E3's. \texttt{train\_sorghum\_4m.py} passes it to the
  train dataset only; a capped val/test would score the arms on
  different yardsticks.
\item
  \texttt{model.model\_size} \(\in\) \{\texttt{small}, \texttt{base},
  \texttt{large}, \texttt{giant}\} (3M) or \{\texttt{small},
  \texttt{base}, \texttt{large}\} (4M). This is what E4 varies: 4M
  builds at 25.1 M / 114.3 M / 332.2 M params.
\item
  \texttt{model.mask\_ratio} is the \textbf{total} masking fraction
  across all modalities; the per-modality split is sampled from
  \texttt{Dirichlet(\$\textbackslash{}alpha\$=dirichlet\_alpha)} once
  per batch.
\item
  \texttt{model.active\_modalities} --- comma-separated subset of
  \texttt{pc,rgb,depth,text}. Restricting it is what the E2 ablation
  varies; everything else stays fixed.
\item
  \texttt{model.loss\_name} \(\in\) \{\texttt{chamfer},
  \texttt{qal\_loss}\} with \texttt{qal\_threshold} /
  \texttt{qal\_alpha} / \texttt{qal\_use\_squared}. Current runs use
  \texttt{qal\_loss}.
\item
  \texttt{model.pc\_loss\_weight} scales the PC term. Chamfer values are
  tiny relative to MSE, so this is typically O(10) when Chamfer is
  selected.
\item
  \texttt{model.spline\_loss\_weight} (4M only) scales the Smooth-L1
  param-regression loss.
\item
  \texttt{model.depth\_norm\_type} \(\in\) \{\texttt{minmax},
  \texttt{standard}\} --- applied to the \textbf{target} before depth
  MSE; the model learns to predict normalised values directly.
\item
  \texttt{checkpointing.resume}: path to a \texttt{.pth} to resume from,
  or \texttt{null} for scratch.
\item
  \texttt{distributed.world\_size\ \textgreater{}\ 1} triggers DDP.
  \texttt{train\_sorghum\_4m.py} also accepts being launched under
  \texttt{torchrun}, in which case \texttt{LOCAL\_RANK} is honoured and
  \texttt{world\_size} is inferred from env.
\end{itemize}

\textbf{Global batch is \texttt{batch\_size\ ×\ world\_size}.}
\texttt{batch\_size} in the YAML is per-GPU, so an ablation must adjust
it to the GPU count to keep the global batch constant --- a global-batch
difference between arms confounds the comparison. The E2 launchers use
16×2 on the 2-GPU Blackwell nodes and 8×4 on the 4-GPU A100 nodes, both
reaching 32.

\textbf{Size an ablation in optimizer steps, not epochs --- and do not
``fix'' epoch counts that disagree.} Under \texttt{view\_sampling} an
epoch is one view per plant, so \emph{epoch size is the plant count}: at
global batch 32 the full train split gives 329 steps/epoch, but a 1
000-plant subset gives 32. Arms that differ in data scale therefore need
very different epoch counts to receive the same number of gradient
updates, and the E3 configs look wrong at a glance because of it:

{\def\LTcaptype{none} % do not increment counter
\begin{longtable}[]{@{}lllll@{}}
\toprule\noalign{}
arm & plants / model & steps/ep & epochs & total steps \\
\midrule\noalign{}
\endhead
\bottomrule\noalign{}
\endlastfoot
\texttt{e3\_1k} & 1 000 & 32 & 6 169 & 197 408 \\
\texttt{e3\_3k} & 3 000 & 94 & 2 100 & 197 400 \\
\texttt{e3\_10k} & 10 000 & 313 & 631 & 197 503 \\
\texttt{e4\_small} / \texttt{e4\_large} & 25.1 M / 332.2 M & 329 & 600 &
197 400 \\
\texttt{e2\_pcrgbdt} & 10 500, base & 329 & 600 & 197 400 \\
\end{longtable}
}

Equalising those epoch counts would hand the 10k arm ten times the
updates of the 1k arm, and the resulting curve would measure data scale
and compute jointly with no way to separate them afterwards.
\texttt{warmup\_epochs}, \texttt{val\_freq} and \texttt{save\_freq} are
scaled by the same per-arm factor, so every arm warms over the same
fraction of its schedule and lands \textasciitilde24 val points.

That shared 197 400-step budget is \texttt{e2\_pcrgbdt}'s, which is why
\textbf{\texttt{e2\_pcrgbdt} is also E3's full-data point and E4's
base-model point} rather than a separate run --- and why the E3/E4
launcher's resources are deliberately byte-identical to the E2 one.
\texttt{e3\_10k} (10 000 plants) and \texttt{e2\_pcrgbdt} (10 500) are
within 5 \% of each other, so that pair doubles as the only run-to-run
error bar in the experiment matrix.

\texttt{persistent\_workers} is deliberately \textbf{off} in the
dataloaders and must stay off: persistent workers hold a copy of the
dataset made at first iteration, so \texttt{train\_ds.set\_epoch(epoch)}
would never reach them and view sampling would silently freeze at epoch
0. The cost is that workers respawn every epoch, which matters most for
\texttt{e3\_1k}, whose epoch is only \textasciitilde1 000 items.

\subsection{Architecture (the part you'd otherwise have to read 4 files
to
learn)}\label{architecture-the-part-youd-otherwise-have-to-read-4-files-to-learn}

\subsubsection{Forward pass shape}\label{forward-pass-shape}

\begin{enumerate}
\def\labelenumi{\arabic{enumi}.}
\tightlist
\item
  \textbf{Embed each modality} to \texttt{(B,\ L\_m,\ D)}:

  \begin{itemize}
  \tightlist
  \item
    RGB / Depth \(\to\) \texttt{PatchEmbed} (Conv2d patchification
    \(\to\) 196 tokens for 224×224, patch=16).
  \item
    Point cloud \(\to\) \texttt{PointCloudEmbed}: FPS samples
    \texttt{num\_pc\_tokens=196} centres, kNN groups
    \texttt{group\_size=32} neighbours per centre, two PointNet-style
    Conv1d blocks produce per-token features.
  \item
    (4M only) Spline params \(\to\) \texttt{TextLeafEmbed}: char
    embeddings + learned positional \(\to\) mean-pool \(\to\) linear
    \(\to\) D. \texttt{n\_text\_tokens\ =\ 1\ +\ max\_leaves} (1 plant
    token + up to 24 leaf tokens).
  \end{itemize}
\item
  \textbf{Add positional + modality embeddings} (each modality has its
  own learned modality bias).
\item
  \textbf{Dirichlet masking} in \texttt{random\_masking\_dirichlet}: a
  single Dirichlet draw decides the visible-token split across
  modalities for the whole batch step. Each sample's visible indices are
  independently random, but the per-modality count is identical
  batch-wide --- this avoids zero-token batch elements that would
  corrupt encoder norms. The 4M variant also enforces
  \texttt{min\_mask\_ratio=0.25} per modality.
\item
  \textbf{Encoder}: visible tokens from all modalities are concatenated
  with a CLS token and fed through \texttt{depth} shared transformer
  blocks (\texttt{embed\_dim=768} for base).
\item
  \textbf{Decoder}: project to \texttt{decoder\_embed\_dim=512}, restore
  mask tokens at original positions per modality, add decoder positional
  embeddings, run \texttt{decoder\_depth=8} transformer blocks, then
  split back into modality-specific heads:

  \begin{itemize}
  \tightlist
  \item
    RGB / Depth \(\to\) linear \(\to\) \texttt{patch\_size²\ ×\ C} per
    token (unpatchified for visualisation).
  \item
    PC \(\to\) FoldingNet-style upsampling: each of 196 PC tokens
    generates
    \texttt{points\_per\_token\ =\ target\_points\ //\ num\_pc\_tokens}
    3D points by concatenating its 512-D feature with a 2D grid
    coordinate and passing through an MLP. \texttt{target\_points}
    defaults to 10000, trimmed/padded to exactly that on output.
  \item
    (4M only) Spline \(\to\) 2-layer MLP \(\to\) \texttt{N\_PARAMS=9}
    floats per token, range {[}0, 1{]}.
  \end{itemize}
\end{enumerate}

\subsubsection{\texorpdfstring{Losses
(\texttt{forward\_loss})}{Losses (forward\_loss)}}\label{losses-forward_loss}

\begin{itemize}
\tightlist
\item
  \textbf{RGB}: per-patch MSE, optionally with \texttt{norm\_pix\_loss}
  (per-patch normalisation), masked mean.
\item
  \textbf{Depth}: per-patch MSE on \textbf{normalised} target (per-image
  min-max or standard); the prediction is compared to the normalised
  target directly.
\item
  \textbf{PC}: bidirectional Chamfer (or QAL) on the full
  \texttt{(B,\ target\_points,\ 3)} cloud (no masking --- the whole
  cloud is reconstructed every step), scaled by
  \texttt{pc\_loss\_weight}.
\item
  \textbf{(4M)}: Smooth-L1 on params, masked by \texttt{text\_valid}
  (real leaf tokens only) \textbf{and} \texttt{mask\_text} (model only
  loses on tokens it didn't see). Scaled by
  \texttt{spline\_loss\_weight}.
\end{itemize}

\texttt{embodied\_mae\_4m.py} imports \texttt{PatchEmbed},
\texttt{PointCloudEmbed}, \texttt{TransformerBlock},
\texttt{chamfer\_distance}, and \texttt{get\_2d\_sincos\_pos\_embed}
from \texttt{embodied\_mae.py} --- when changing these, expect both
models to be affected.

\subsubsection{Param normalisation (4M
only)}\label{param-normalisation-4m-only}

Plant and leaf parameters are normalised to {[}0, 1{]} via fixed
\texttt{\_PLANT\_SCALE\ /\ \_PLANT\_SHIFT} and
\texttt{\_LEAF\_SCALE\ /\ \_LEAF\_SHIFT} arrays at the top of
\texttt{embodied\_mae\_4m.py}. \texttt{decode\_params\_to\_text}
un-normalises them back into the human-readable strings shown in
visualisations. If new fields are added to the spline YAMLs, both arrays
and the \texttt{\_plant\_to\_params} / \texttt{\_leaf\_to\_params}
builders must be updated.

\subsection{Dataset layout}\label{dataset-layout}

\texttt{SorghumDataset} expects:

\begin{verbatim}
data_root/{train,val,test}/<sample_name>/
    rgb.png              # 224-ready RGB render
    depth.png            # big-endian packed RGBA depth
    *_nc_cam.ply         # camera-frame, normals-cleaned point cloud
    *_spline.yml         # 4M only — procedural generation params
\end{verbatim}

Sample folders are named
\texttt{Sorghum\_\textless{}plant\textgreater{}\_\textless{}view\textgreater{}},
so \texttt{Sorghum\_0\_00\ …\ Sorghum\_0\_09} are ten views of one
plant. \textbf{Only those four files are read by the loaders.} The
\texttt{Sorghum\_\textless{}n\textgreater{}.obj} and
\texttt{Sorghum\_\textless{}n\textgreater{}\_nc.ply} that sit alongside
them are source assets, are duplicated in full into every one of a
plant's ten view folders, and are never opened during training --- they
are \textasciitilde64 \% of the bytes on disk.

The PC loader uniformly samples / pads to \texttt{num\_points}, then
centres at the centroid and scales so the max-distance point lands on
the unit sphere. RGB uses ImageNet mean/std normalisation; depth is
loaded as single-channel L and only \texttt{ToTensor}'d (no
normalisation at load --- normalisation happens inside the loss).

The active dataset on Nova is
\texttt{/work/mech-ai-scratch/alloy/shorgum\_data/new\_data\_50K/Sorghum\_15K/},
an extreme-enriched 70/15/15 split (seed 42) of 15 000 plants:
\textbf{105 000 train / 22 500 val / 22 500 test} view-samples.
\texttt{assignment.csv} and \texttt{features.csv} at that root record
the split; the tooling to regenerate or reshuffle it is in
\texttt{data\_split/}.

\subsection{Porting to another
machine}\label{porting-to-another-machine}

Everything machine-specific is in three places --- nothing else needs
touching:

\begin{enumerate}
\def\labelenumi{\arabic{enumi}.}
\tightlist
\item
  \textbf{\texttt{data.data\_root} in the YAML you run.} Every config
  hardcodes the Nova absolute path above.
\item
  \textbf{The \texttt{\#SBATCH} headers in \texttt{slurm/*.sbatch}} ---
  \texttt{-\/-partition}, \texttt{-\/-account}, \texttt{-\/-gres},
  \texttt{-\/-cpus-per-task}, \texttt{-\/-mem}. These encode Nova's
  partitions (\texttt{nova}, \texttt{scavenger}) and accounts
  (\texttt{mech-ai}, \texttt{mech-ai-scavenger}) and mean nothing
  elsewhere. On a non-SLURM box, ignore \texttt{slurm/} entirely and use
  the \texttt{torchrun} line above.
\item
  \textbf{The conda env.} \texttt{environment.yml} rebuilds
  \texttt{det}; it pins a CUDA 12.4 PyTorch, so a different GPU
  generation may need a different build (see the sm\_120 note under
  Environment).
\end{enumerate}

\textbf{Moving the data is the expensive part.} At \textasciitilde5.0 MB
per sample folder × 150 000 folders the split is
\textbf{\textasciitilde750 GB}, but the four files training actually
reads total \textasciitilde1.8 MB per sample, so a filtered copy is
\textbf{\textasciitilde265 GB} --- under 40 \% of the naive transfer.
Copy with an include-filter rather than syncing the tree:

\begin{Shaded}
\begin{Highlighting}[]
\FunctionTok{rsync} \AttributeTok{{-}a} \AttributeTok{{-}{-}info}\OperatorTok{=}\NormalTok{progress2 }\DataTypeTok{\textbackslash{}}
  \AttributeTok{{-}{-}include}\OperatorTok{=}\StringTok{\textquotesingle{}*/\textquotesingle{}} \DataTypeTok{\textbackslash{}}
  \AttributeTok{{-}{-}include}\OperatorTok{=}\StringTok{\textquotesingle{}rgb.png\textquotesingle{}} \AttributeTok{{-}{-}include}\OperatorTok{=}\StringTok{\textquotesingle{}depth.png\textquotesingle{}} \DataTypeTok{\textbackslash{}}
  \AttributeTok{{-}{-}include}\OperatorTok{=}\StringTok{\textquotesingle{}*\_nc\_cam.ply\textquotesingle{}} \AttributeTok{{-}{-}include}\OperatorTok{=}\StringTok{\textquotesingle{}*\_spline.yml\textquotesingle{}} \DataTypeTok{\textbackslash{}}
  \AttributeTok{{-}{-}exclude}\OperatorTok{=}\StringTok{\textquotesingle{}*\textquotesingle{}} \DataTypeTok{\textbackslash{}}
\NormalTok{  /path/to/Sorghum\_15K/ user@host:/dest/Sorghum\_15K/}
\end{Highlighting}
\end{Shaded}

Also copy \texttt{assignment.csv} and \texttt{features.csv} from the
split root, and any warm-start checkpoint you need
(\texttt{best\_model.pth} is \textasciitilde1.3 GB per run).

Note the dataset loaders build a folder index on first use and cache it
--- the first run on a fresh copy pays a one-off scan (\textasciitilde24
min for the 105 k train split); later runs print
\texttt{index\ cache\ hit}. The pipeline is \textbf{dataloader-bound,
not GPU-bound}: \texttt{num\_workers} (\$\approx\$1.6 items/s per
worker) sets throughput, so give it as many CPUs as the node allows and
scale \texttt{-\/-mem} with the worker count (\textasciitilde2.3 GB RSS
per worker plus \textasciitilde10 GB per node for model and CUDA
context).

\subsection{Outputs}\label{outputs}

Each run writes to \texttt{\textless{}output\_dir\textgreater{}/}: -
\texttt{checkpoints/checkpoint\_epoch\_\textless{}N\textgreater{}.pth}
every \texttt{save\_freq} epochs - \texttt{best\_model.pth} when val
loss improves -
\texttt{visualizations/epoch\_\textless{}N\textgreater{}\_sample\_\textless{}i\textgreater{}\_\textless{}name\textgreater{}.png}
every \texttt{viz\_freq} epochs (4-row grid for 3M, 5-row grid for 4M
including text predictions); skipped automatically when the run is a
reduced-modality arm - \texttt{training\_history.json} (rolling) -
\texttt{config.json} --- snapshot of the effective args. \textbf{Check
this to confirm what a run actually used}, especially
\texttt{batch\_size\ ×\ world\_size}, \texttt{active\_modalities},
\texttt{max\_plants} and \texttt{model\_size}.

\texttt{outputs/} is gitignored apart from a small whitelist in
\texttt{.gitignore}. Wandb logging is on by default
(\texttt{use\_wandb:\ true}, project \texttt{embodied-mae-sorghum});
project and run names differ between runs --- check the YAML, not the
script defaults.

\subsection{Things that look like dead code but
aren't}\label{things-that-look-like-dead-code-but-arent}

\begin{itemize}
\tightlist
\item
  \texttt{outputs\_sorghum\_*/} directories at the repo root are old run
  outputs kept for reference; the canonical output root is
  \texttt{./outputs/}.
\item
  \texttt{tools/process\_depth\_bg.py}, \texttt{tools/process\_mask.py},
  \texttt{tools/validate\_sorghum\_data.py}, \texttt{sweeps/vis.py},
  \texttt{tools/check\_structure.py} are one-off data-prep / diagnostic
  scripts, not part of any pipeline.
\item
  \texttt{sweeps/vis\_pc\_masking.py} is a standalone tool for
  visualising the FPS + Dirichlet masking on a single point cloud.
\item
  \texttt{tools/visualize\_sorghum\_pointclouds.py} renders multi-view
  PC galleries from raw \texttt{.ply} files; it doesn't touch the model.
\item
  \texttt{eval/analyze\_e2.py} reads the E2 arm output dirs and builds
  the modality-value-add comparison.
\end{itemize}

\end{document}
